% GENERATED FILE. Edit canonical lessons, then run npm run generate:books.
% canonical-source-hash: fnv1a64:682a5bf29b2b2858
% canonical-lessons: HI-C540-tajurba, HI-C540-yukti, HI-C540-zariya, HI-C540-tarkib, HI-C540-prayatn

\chapter{Ways and Means}
\label{ch:hi-ways-and-means}
\hlchaptermodality{540}

\begin{chapteropening}
\textbf{By the end of this chapter:} \emph{I can say how I tried to get something done.}

Complete the last lesson of chapter 540: I can say how I tried to get something done.
\end{chapteropening}

\section[\texorpdfstring{\hi{तजुर्बा} (tajurbā)}{तजुर्बा (tajurbā)}]{\hi{तजुर्बा} (tajurbā) — experience (gained over time)}

\label{lesson:HI-C540-tajurba}



\begin{quote}
Before the new one: say the Hindi for a wish (for someone's good), then the Hindi for to resolve.
\end{quote}

\subsection*{You'll want to know: \hi{तजुर्बा}}
\textbf{\hi{तजुर्बा}} — \emph{tajurbā} — \textquotedblleft{}experience (gained over time)\textquotedblright{}.

\begin{quote}
\hi{अरुण} \hi{को} \hi{इस} \hi{काम} \hi{का} \hi{दस} \hi{साल} \hi{का} \hi{तजुर्बा} \hi{है।}
\end{quote}

\begin{quote}
\emph{Aruṇ ko is kām kā das sāl kā tajurbā hai.}
\end{quote}

\begin{quote}
Arun has ten years' experience of this work.
\end{quote}

Masculine.

\begin{cousinweb}
\textbf{\hi{तजुर्बा}} (\emph{tajurbā}) is usually traced to Arabic \emph{tajriba}, \textquotedblleft{}a trial, a test\textquotedblright{}, through Persian; Hindi uses it for what years of doing something teach you.
\end{cousinweb}

\begin{grammarlens}[title={What you've built}]
One more word for saying what happens if.
\end{grammarlens}

\subsection*{Your turn}
\begin{itemize}
\raggedright
  \item \emph{Say it:} \emph{tajurbā}
  \item \emph{Say it:} \emph{tajurbā}, once more
  \item \emph{Say it:} \emph{Aruṇ ko is kām kā das sāl kā tajurbā hai.} — Arun has ten years' experience of this work.
  \item \emph{Recall:} say the Hindi for a wish (for someone's good), then the Hindi for to resolve, then say \emph{tajurbā} again
\end{itemize}

\begin{tcolorbox}[breakable,colback=teal!4,colframe=teal!35!black,title={Before you move on}]
What does \hi{तजुर्बा} mean? (\textquotedblleft{}Experience (gained over time)\textquotedblright{}.) Say it once more.
\end{tcolorbox}
\section[\texorpdfstring{\hi{युक्ति} (yukti)}{युक्ति (yukti)}]{\hi{युक्ति} (yukti) — a clever way; a device}

\label{lesson:HI-C540-yukti}



\begin{quote}
Before the new one: say the Hindi for to resolve, then the Hindi for experience (gained over time).
\end{quote}

\subsection*{You'll want to know: \hi{युक्ति}}
\textbf{\hi{युक्ति}} — \emph{yukti} — \textquotedblleft{}a clever way; a device\textquotedblright{}.

\begin{quote}
\hi{अरुण} \hi{ने} \hi{एक} \hi{युक्ति} \hi{से} \hi{भारी} \hi{मेज़} \hi{अकेले} \hi{उठाई।}
\end{quote}

\begin{quote}
\emph{Aruṇ ne ek yukti se bhārī mez akele uṭhāī.}
\end{quote}

\begin{quote}
Arun found a clever way to lift the heavy table on his own.
\end{quote}

Feminine.

\begin{cousinweb}
\textbf{\hi{युक्ति}} (\emph{yukti}) is a tatsama, taken whole from Sanskrit \emph{yukti}, \textquotedblleft{}union; a contrivance, a means\textquotedblright{} (Monier-Williams).
\end{cousinweb}

\begin{grammarlens}[title={What you've built}]
One more word for saying what happens if.
\end{grammarlens}

\subsection*{Your turn}
\begin{itemize}
\raggedright
  \item \emph{Say it:} \emph{yukti}
  \item \emph{Say it:} \emph{yukti}, once more
  \item \emph{Say it:} \emph{Aruṇ ne ek yukti se bhārī mez akele uṭhāī.} — Arun found a clever way to lift the heavy table on his own.
  \item \emph{Recall:} say the Hindi for to resolve, then the Hindi for experience (gained over time), then say \emph{yukti} again
\end{itemize}

\begin{tcolorbox}[breakable,colback=teal!4,colframe=teal!35!black,title={Before you move on}]
What does \hi{युक्ति} mean? (\textquotedblleft{}A clever way; a device\textquotedblright{}.) Say it once more.
\end{tcolorbox}
\section[\texorpdfstring{\hi{ज़रिया} (zariyā)}{ज़रिया (zariyā)}]{\hi{ज़रिया} (zariyā) — a means, a way (of doing something)}

\label{lesson:HI-C540-zariya}



\begin{quote}
Before the new one: say the Hindi for experience (gained over time), then the Hindi for a clever way.
\end{quote}

\subsection*{You'll want to know: \hi{ज़रिया}}
\textbf{\hi{ज़रिया}} — \emph{zariyā} — \textquotedblleft{}a means, a way (of doing something)\textquotedblright{}.

\begin{quote}
\hi{मोबाइल} \hi{बात} \hi{करने} \hi{का} \hi{सबसे} \hi{आसान} \hi{ज़रिया} \hi{है।}
\end{quote}

\begin{quote}
\emph{Mobāil bāt karne kā sabse āsān zariyā hai.}
\end{quote}

\begin{quote}
A mobile is the easiest way to talk.
\end{quote}

Masculine.

\begin{cousinweb}
\textbf{\hi{ज़रिया}} (\emph{zariyā}) is usually traced to Arabic \emph{ḏarīʿa}, \textquotedblleft{}a means\textquotedblright{}, through Persian.
\end{cousinweb}

\begin{grammarlens}[title={What you've built}]
One more word for saying what happens if.
\end{grammarlens}

\subsection*{Your turn}
\begin{itemize}
\raggedright
  \item \emph{Say it:} \emph{zariyā}
  \item \emph{Say it:} \emph{zariyā}, once more
  \item \emph{Say it:} \emph{Mobāil bāt karne kā sabse āsān zariyā hai.} — A mobile is the easiest way to talk.
  \item \emph{Recall:} say the Hindi for experience (gained over time), then the Hindi for a clever way, then say \emph{zariyā} again
\end{itemize}

\begin{tcolorbox}[breakable,colback=teal!4,colframe=teal!35!black,title={Before you move on}]
What does \hi{ज़रिया} mean? (\textquotedblleft{}A means, a way (of doing something)\textquotedblright{}.) Say it once more.
\end{tcolorbox}
\section[\texorpdfstring{\hi{तरकीब} (tarkīb)}{तरकीब (tarkīb)}]{\hi{तरकीब} (tarkīb) — a trick; a method (for doing something)}

\label{lesson:HI-C540-tarkib}



\begin{quote}
Before the new one: say the Hindi for a clever way, then the Hindi for a means, a way (of doing something).
\end{quote}

\subsection*{You'll want to know: \hi{तरकीब}}
\textbf{\hi{तरकीब}} — \emph{tarkīb} — \textquotedblleft{}a trick; a method (for doing something)\textquotedblright{}.

\begin{quote}
\hi{मीरा} \hi{के} \hi{पास} \hi{हर} \hi{समस्या} \hi{की} \hi{कोई} \hi{तरकीब} \hi{होती} \hi{है।}
\end{quote}

\begin{quote}
\emph{Mīrā ke pās har samasyā kī koī tarkīb hotī hai.}
\end{quote}

\begin{quote}
Mira always has some trick for every problem.
\end{quote}

Feminine.

\begin{cousinweb}
\textbf{\hi{तरकीब}} (\emph{tarkīb}) is usually traced to Arabic \emph{tarkīb}, \textquotedblleft{}putting together\textquotedblright{}, through Persian.
\end{cousinweb}

\begin{grammarlens}[title={What you've built}]
One more word for saying what happens if.
\end{grammarlens}

\subsection*{Your turn}
\begin{itemize}
\raggedright
  \item \emph{Say it:} \emph{tarkīb}
  \item \emph{Say it:} \emph{tarkīb}, once more
  \item \emph{Say it:} \emph{Mīrā ke pās har samasyā kī koī tarkīb hotī hai.} — Mira always has some trick for every problem.
  \item \emph{Recall:} say the Hindi for a clever way, then the Hindi for a means, a way (of doing something), then say \emph{tarkīb} again
\end{itemize}

\begin{tcolorbox}[breakable,colback=teal!4,colframe=teal!35!black,title={Before you move on}]
What does \hi{तरकीब} mean? (\textquotedblleft{}A trick; a method (for doing something)\textquotedblright{}.) Say it once more.
\end{tcolorbox}
\section[\texorpdfstring{\hi{प्रयत्न} (prayatn)}{प्रयत्न (prayatn)}]{\hi{प्रयत्न} (prayatn) — an effort, an endeavour (formal)}

\label{lesson:HI-C540-prayatn}



\begin{quote}
Before the new one: say the Hindi for a means, a way (of doing something), then the Hindi for a trick.
\end{quote}

\subsection*{You'll want to know: \hi{प्रयत्न}}
\textbf{\hi{प्रयत्न}} — \emph{prayatn} — \textquotedblleft{}an effort, an endeavour (formal)\textquotedblright{}.

\begin{quote}
\hi{बहुत} \hi{प्रयत्न} \hi{के} \hi{बाद} \hi{वह} \hi{पहाड़} \hi{पर} \hi{पहुँचा।}
\end{quote}

\begin{quote}
\emph{Bahut prayatn ke bād vah pahāṛ par pahũcā.}
\end{quote}

\begin{quote}
After a great effort he reached the top of the mountain.
\end{quote}

Masculine.

A formal word.

\begin{cousinweb}
\textbf{\hi{प्रयत्न}} (\emph{prayatn}) is a tatsama, taken whole from Sanskrit \emph{prayatna}, \textquotedblleft{}persevering effort\textquotedblright{} (Monier-Williams).
\end{cousinweb}

\begin{grammarlens}[title={What you've built}]
One more word for saying what happens if.
\end{grammarlens}

\subsection*{Your turn}
\begin{itemize}
\raggedright
  \item \emph{Say it:} \emph{prayatn}
  \item \emph{Say it:} \emph{prayatn}, once more
  \item \emph{Say it:} \emph{Bahut prayatn ke bād vah pahāṛ par pahũcā.} — After a great effort he reached the top of the mountain.
  \item \emph{Recall:} say the Hindi for a means, a way (of doing something), then the Hindi for a trick, then say \emph{prayatn} again
\end{itemize}

\begin{tcolorbox}[breakable,colback=teal!4,colframe=teal!35!black,title={Before you move on}]
What does \hi{प्रयत्न} mean? (\textquotedblleft{}An effort, an endeavour (formal)\textquotedblright{}.) Say it once more.
\end{tcolorbox}
